% =============================================================================
% APPENDIX — Reference Artifacts
% Every prompt and schema used in the study, reproduced verbatim for
% transparency. Plus the real-data trace of one vignette.
% =============================================================================

\subsection{System Prompts}
\label{app:system_prompts}

The system prompts below differentiate conditions B$'$, C, and D. Conditions
A, B, and E use no system prompt (the dataset's per-row
\emph{confinement instruction}, described in \cref{sec:methods_benchmark},
is the only prompt suffix appended).

\begin{lstlisting}[caption={Condition B$'$ --- force tool use. Verbatim
from \toolname{prompts/condition\_b\_prime.txt}.},
label={lst:prompt_bprime}, language=]
You are an expert clinician answering medical calculator questions for
Indian patients.

You must call the appropriate calculator tool for every computation — do
not compute values by hand, even if you believe the answer is obvious.
Always invoke the matching `medical_calculator_*` tool with inputs taken
directly from the patient case, and use the tool's returned value as
your answer.

Read the case carefully, and after the tool returns, reply strictly
with the JSON object the user requests.
\end{lstlisting}

\begin{lstlisting}[caption={Condition C --- prompt-only self-verify.
Verbatim from \toolname{prompts/condition\_c.txt}.},
label={lst:prompt_c}, language=]
You are an expert clinician answering medical calculator questions for
Indian patients.

Before calling any calculator tool, verify each input you plan to pass
against the patient case: confirm the units match (e.g. mg/dL vs
mmol/L), the enum choices are correct (e.g. sex, activity level,
severity), and the numeric values are taken directly from the vignette.
If the case lacks clear evidence for a required input, do not guess —
re-read the case before calling the tool.

Read the case carefully, use any available tools if useful, and reply
strictly with the JSON object the user requests.
\end{lstlisting}

\begin{lstlisting}[caption={Condition D --- post-hoc Sonnet verifier
system prompt. Verbatim from \toolname{prompts/condition\_d.txt}.},
label={lst:prompt_d}, language=]
You are a verification reviewer for clinical calculator answers.

You will be shown (1) a patient case in Indian clinical language, (2)
the question's required output format, and (3) a candidate JSON answer
produced by another model. Your job is to identify whether the
candidate answer is consistent with the case — specifically, whether
the inputs the answering model used (if visible in its reasoning) match
what the case states.

Reply with a JSON object of the form:
{
  "consistent": true | false,
  "issue": "<one-sentence description of the discrepancy if
   consistent=false, otherwise empty string>"
}

Be conservative: only flag clear inconsistencies in units, enum choices
(sex, activity level, severity, etc.), or numeric values that
contradict the case. Do not flag minor formatting differences or
stylistic choices. Do not flag a missing field unless the question
requires it. If you are not certain there is an inconsistency, return
consistent=true.
\end{lstlisting}

\subsection{Fact Extractor Prompt}
\label{app:extractor_prompt}

The fact extractor used in E, B$'$+E, and B$'$+E (reactive) runs Sonnet
with the system prompt shown in \cref{lst:prompt_extractor_header},
followed by a schema body that is generated deterministically from the
MCP per-calculator JSON Schemas. For the upfront variant the body is
the full union of all calculator fields ($\sim$500 entries); for the
reactive variant it is the subset of fields the active tool call
requests (typically 3--30). The header and rules are byte-identical
across the two variants; only the schema body differs.

\begin{lstlisting}[caption={Fact extractor system prompt header
(stable across upfront and reactive variants). The schema body that
follows is generated programmatically from the MCP schema dump by
\toolname{harness/extraction/prompt\_builder.py}.},
label={lst:prompt_extractor_header}, language=]
You are a medical fact extractor. Read the patient vignette and return
a strict JSON object capturing every value the case explicitly states
for the fields below.

Return ONLY a JSON object (no prose, no markdown fences). Include ONLY
the fields the case explicitly states a value for. Omit any field the
case does not state — do not emit null for absent fields. The schema
below is the vocabulary of allowed field names; you do not need to
emit every field.

Rules:
- Include only what the case explicitly states. If a value must be
  inferred, computed, or guessed, omit the field.
- Use the units the case provides. Do not convert units. If the case
  says "creatinine 130 µmol/L", store 130, not the converted mg/dL value.
- For ranges (e.g. "BP 140/90"), split into the appropriate fields.
- For descriptors that map to a calculator's enum (e.g. "sedentary
  lifestyle" → "sedentary"), use the enum value if the mapping is
  unambiguous; otherwise omit.
- If the case uses Hinglish or non-standard phrasing, interpret in
  standard clinical context where unambiguous; otherwise omit.
- Return only valid JSON. No commentary.

Schema (allowed field names with type/range/enum constraints):
{ ...generated schema body... }
\end{lstlisting}

\subsection{Verifier Feedback Template (Condition E)}
\label{app:feedback_template}

The verifier's flag injects the following template into the primary
model's prompt. The \toolname{\{violations\}} placeholder expands to a
list of disagreeing fields, each annotated with the planned value, the
extracted value, and a short diagnostic note.

\begin{lstlisting}[caption={Verifier feedback prompt, verbatim from
\toolname{prompts/condition\_e\_feedback.txt}.},
label={lst:feedback}, language=]
Your planned tool call has an input that does not match the patient case:

{violations}

Re-read the case carefully and call the tool again with corrected
inputs. If the case truly does not specify the required input, do not
invent a value; explain in your reply.
\end{lstlisting}

\subsection{Example MCP Calculator Schema}
\label{app:bmi_schema}

Across all $\sim$70 calculators in MedAI MCP, the union of property
names exceeds 500 unique fields; the schema for any single calculator
typically declares 3--30. One representative calculator's JSON Schema
is reproduced verbatim below.

\begin{lstlisting}[caption={Abbreviated MCP JSON Schema for
\toolname{bmi\_calculator\_body\_mass\_index}. The full schema includes
descriptions, units, and range constraints for each field. Note that
\texttt{sex} is declared but \emph{not} required; the BMI calculator
will dispatch without it.}, label={lst:bmi_schema}, language=]
{
  "name": "bmi_calculator_body_mass_index",
  "input_data": {
    "type": "object",
    "required": ["height", "mass"],
    "properties": {
      "sex":     { "$ref": "#/$defs/Sex", "description": "Biological sex" },
      "height":  { "type": "number",  "description": "Height in meters" },
      "mass":    { "type": "number",  "description": "Body weight in kilograms" },
      "age":     { "type": "integer", "description": "Age in years" },
      "athlete": { "$ref": "#/$defs/YesNo", "description": "Athletic status" }
    }
  }
}
\end{lstlisting}

\subsection{Real-Data Trace: Vignette \texttt{bmi\_001}}
\label{app:loop_trace}

% Real-data sequence trace for §5.3 using vignette bmi_001.
\begin{figure}[H]
\centering
\begin{tikzpicture}[
  font=\sffamily\footnotesize,
  >={Stealth[length=2.0mm]},
  every node/.style={align=left},
  node distance=3.5mm,
  stepbox/.style={rectangle, rounded corners=2pt, line width=0.5pt,
               inner sep=4.5pt, text width=118mm},
  case/.style={stepbox, draw=warmGray, fill=paleGray},
  primary/.style={stepbox, draw=accent, fill=accentLight},
  verifier/.style={stepbox, draw=condBprimeE, fill=condBprimeE!10},
  bad/.style={stepbox, draw=condBprimeE, fill=condBprimeE!18, line width=0.8pt},
  good/.style={stepbox, draw=condReactive, fill=condReactive!12},
]

\node[case] (case) {%
  \textbf{Vignette \texttt{bmi\_001}} (real exemplar from the bundle).\\
  \textit{Case:} A 29-year-old \textbf{female} patient, weight 68\,kg,
  height 1.6\,m. Compute her BMI.\\
  \textit{Ground truth:} BMI $= 26.6$, tolerance $\pm 0.1$.};

\node[primary, below=of case] (s1) {%
  \textbf{Step 1.}\ \ Upfront extractor (Sonnet, full $\sim$500-field
  schema) populates a reference set including \texttt{sex:\,female,
  age:\,29, weight:\,68, height:\,1.6, bmi:\,26.56}. Sex is extracted
  correctly.};

\node[primary, below=of s1] (s2) {%
  \textbf{Step 2.}\ \ Qwen3 (under B$'$) emits a tool call to
  \toolname{medical\_calculator\_output} with
  \texttt{input\_data=\{mass:\,68, height:\,1.6, sex:\,1\}}. The model
  encodes \texttt{sex} as an integer (its preferred binary encoding)
  even though the BMI calculator does not require this field.};

\node[verifier, below=of s2] (s3) {%
  \textbf{Step 3.}\ \ Verifier flags: extracted \texttt{sex=`female'}
  (string) does not match planned \texttt{sex=1} (int). Note: ``value
  type differs from the case (str vs int).'' MCP call is \emph{not}
  dispatched; feedback re-prompt is injected.};

\node[verifier, below=of s3] (s4) {%
  \textbf{Step 4.}\ \ The cycle repeats eight times. The model tries
  alternate encodings, the verifier keeps flagging on type mismatch
  against the extractor's string reference, and the turn budget exhausts
  before the BMI calculator ever runs.};

\node[bad, below=of s4] (s5) {%
  \textbf{Outcome (B$'$+E):}\ \ Final answer \texttt{nan}.
  \emph{This vignette was correct under B$'$ alone}; the verifier
  flipped a correct call into a failure mode.};

\node[good, below=of s5, yshift=-1mm] (s6) {%
  \textbf{Under reactive B$'$+E on the same vignette.}\ \ The model's
  first tool call under reactive carries
  \texttt{input\_data=\{mass:\,68, height:\,1.6, age:\,29\}}; the
  focused extractor is asked only for those three fields, never returns
  a value for \texttt{sex}, and the verifier's ``no extractor value
  $\rightarrow$ no flag'' rule applies. The call dispatches; the BMI
  calculator returns $26.56$. \textit{Final answer: 26.6 (correct).}};

\draw[->, accent]      (case) -- (s1);
\draw[->, accent]      (s1)   -- (s2);
\draw[->, condBprimeE] (s2)   -- (s3);
\draw[->, condBprimeE, dashed] (s3) -- (s4);
\draw[->, condBprimeE] (s4)   -- (s5);

\end{tikzpicture}
\caption{Trace of vignette \texttt{bmi\_001}: an unrequested
\texttt{sex} value in the upfront reference triggers a spurious
flag-loop that reactive extraction avoids.}
\label{fig:loop_trace}
\end{figure}


\Cref{fig:loop_trace} walks through one real exemplar. Under upfront
B$'$+E, the extractor populated a value for \texttt{sex} that the BMI
calculator did not request; the model's tool call used a different
serialization (\texttt{sex=1} vs the extracted \texttt{sex=`female'}),
triggering an unresolvable type-mismatch loop that exhausted the turn
budget. Under reactive extraction the focused prompt for the BMI
calculator did not ask the extractor about \texttt{sex} at all, so no
flag fired and the calculator dispatched cleanly. This is one of 67
vignettes following the recover-by-reactive pattern in our run
(\cref{fig:trajectory_flow}).
